\documentclass[aspectratio=43]{beamer}

\usepackage[T1]{fontenc}
\usepackage[utf8]{inputenc}
\usepackage[english]{babel}
\usepackage{lmodern}
\usepackage{booktabs}
\usepackage{multirow}
\usepackage{subcaption}

%import the theme and packages for the code blocks
\input{code_blocks_theme.tex}

% Use Unipd as theme, with options:
% - pageofpages: define the separation symbol of the footer page of pages (e.g.: of, di, /, default: of)
% - logo: position another logo near the Unipd logo in the title page (e.g. department logo), passing the second logo path as option 
% Use the environment lastframe to add the endframe text
\usetheme[pageofpages=of]{Unipd}

\title{A Comparative Study of Image Inpainting Methods
	under Diverse Masking Conditions\\}
\subtitle{}
\author[]{Ponchio Mattia, Secco Benedetto}

\date{16th September 2026}


\begin{document}
	\frame{\titlepage}
	
	
	\section{Introduction}
	\begin{frame}{Image Inpainting}
		\centering
		\begin{itemize}
			\item Image inpainting is the task of filling in holes in an image.
			\item In non-blind inpainting, the location of the corrupted pixels is provided in order to fill them with plausible imagery.
			\item Useful to edit out unwanted objects or generate occluded regions.
			\item Task has been tackled both with traditional methods (e.g. Biharmonic Equations, PatchMatch) and deep learning networks.
		\end{itemize}
	\end{frame}
	
	
	\begin{frame}{Dataset}
		\centering
		\begin{itemize}
			\item We utilize a subset of the Places365 test set, which contains unique imagery related to 365 scene categories.
			\item In particular, we consider 200 images from the resized test set, where the images have been resized to $256 \times 256$ regardless of the original aspect ratio.
			\item For \textit{each} image, we consider its combination with each one of our 48 masks, totaling 9600 unique inputs to the networks
		\end{itemize}
	\end{frame}
	
	
	\begin{frame}{Masks}
		\begin{columns}[c]
			
			\begin{column}{0.5\textwidth}
				\centering
				\includegraphics[width=\textwidth]{images/grid_output.pdf}
			\end{column}
			
			\begin{column}{0.5\textwidth}
				\begin{itemize}
					\item Masks used with designated index in the top right (masked region in white).
					\item In order sap, box, circle and random walk generation methods.
					\item For each type, filters with masked pixels to total pixels ratio in $\{0.10, 0.25, 0.40\}$ are generated, first not allowing to cross image borders and then permitting it.
				\end{itemize}
			\end{column}
			
		\end{columns}
	\end{frame}
	
	% ===========================================================
	% =============== MASK DESIGN ==============================
	% ==========================================================
	\begin{frame}{Mask Design Choice}
		
		We evaluate models using masks varying in \textbf{type}, \textbf{corruption ratio}, and \textbf{position} (center vs border).
		
		\vspace{0.2cm}
		
		\begin{itemize}
			\item \textbf{SAP (salt-and-pepper):} sparse, independent pixels $\rightarrow$ tests local texture / denoising
			
			\item \textbf{Box / Circle:} large contiguous regions $\rightarrow$ tests structural + semantic completion and the effect of boundary geometry (sharp vs smooth)
			
			\item \textbf{Random Walk:} irregular, thin connected structures $\rightarrow$ tests long-range context and robustness
		\end{itemize}
		
		\vspace{0.2cm}
		
		\textbf{Corruption ratios:} $\{0.10,\;0.25,\;0.40\}$ \hfill
		\textbf{Position:} center vs border
		
	\end{frame}
	
	
	
	\section{Partial Convolution}
	\begin{frame}
		\begin{center}
			\Huge{\textbf{Partial Convolution}}
		\end{center}
	\end{frame}
	
	
	\begin{frame}{Goals}
		\begin{itemize}
			\item Independent from hole initialization values.
			\item No expensive post-processing needed.
			\item Convolution results depend only on the non-hole regions at every layer, with the masked region being automatically updated at each step.
			\item Handle both regular and irregular shaped holes, regardless of their proximity to the borders of the image.
			\item We implement a pre-trained model, which was trained on the Places dataset for 500000 iterations.
		\end{itemize}
	\end{frame}
	
	
	\begin{frame}{Partial Convolutional Layer}
		\begin{center}			
			\begin{equation*}
				x' = \begin{cases}
					\mathbf{W}^T(\mathbf{X}\odot\mathbf{M})\frac{\text{sum}(\mathbf{1})}{\text{sum}(\mathbf{M})} + b,  & \text{if }  \text{sum}(\mathbf{M})>0  \\
					0,  & \text{otherwise}
				\end{cases}
			\end{equation*}
			\begin{itemize}
				\item Mask convention: 0 for masked portion.
				\item Output values depend only on unmasked inputs.
				\item Scaling factor to adjust for the varying amount of unmasked inputs.
				\item If the convolution was able to condition its output on at least one valid input value, then that location becomes valid, expressed with:
			\end{itemize}
			\begin{equation*}
					m' = \begin{cases}
					1,  & \text{if }  \text{sum}(\mathbf{M})>0  \\
					0,  & \text{otherwise}
				\end{cases}
			\end{equation*}
		\end{center}
	\end{frame}
	
	
	\begin{frame}{Network Architecture Inspiration}
		\centering
		\includegraphics[scale=0.14]{images/u-net-architecture.png}
		\begin{itemize}
			\item UNet-like architecture: Encoder-Decoder with skip connections.
		\end{itemize}
		\small
		\footnote{Image source: Ronneberger, O., Fischer, P., Brox, T.: U-net: Convolutional networks for biomedical image segmentation. Springer (2015)}
	\end{frame}
	
	
	\begin{frame}{Network Architecture Details}
		\begin{itemize}
			\item Each convolution is transformed into a \textbf{partial} convolution.
			\item The skip links concatenate two feature maps and two masks respectively, acting as the feature and mask inputs for the next partial convolution layer.
			\item The last partial convolution layer’s input contains the concatenation of the original input image with hole and original mask, making it possible for the model to copy non-hole pixels.
		\end{itemize}
	\end{frame}
	
	
	%\begin{frame}{Network Architecture Visualization}
	%	\centering
	%	\includegraphics[width=0.9\textwidth, height=0.85\textheight]{images/Pconv_net.png}
	%\end{frame}
	
	
	\begin{frame}{Loss Function}
		\centering
		\begin{equation*}
				\mathcal{L}_{total} = 
				\mathcal{L}_{valid} + 6\mathcal{L}_{hole} + 0.05\mathcal{L}_{perceptual}
				+ 120(\mathcal{L}_{style_{out}} + \mathcal{L}_{style_{comp}}) + 0.1\mathcal{L}_{tv}
		\end{equation*}
		\begin{itemize}
			\item Both pixel-wise and perceptual terms.
			\item $\mathcal{L}_{hole}$ and $\mathcal{L}_{valid}$ are the $L_1$ losses on the network output for the hole and the non-hole pixels respectively.
		\begin{equation*}
			\centering
			\mathcal{L}_{perceptual} = \sum_{p=0}^{P-1} \frac{\| \Psi_{p}^{\mathbf{I}_{out}} - \Psi_{p}^{\mathbf{I}_{gt}} \|_1}{N_{\Psi_{p}^{\mathbf{I}_{gt}}}}
			+ \sum_{p=0}^{P-1} \frac{\|  \Psi_{p}^{\mathbf{I}_{comp}} - \Psi_{p}^{\mathbf{I}_{gt}} \|_1}{N_{\Psi_{p}^{\mathbf{I}_{gt}}}}
		\end{equation*}
			\item $\mathbf{I}_{comp}$ is the raw output image $\mathbf{I}_{out}$, but with the non-hole pixels directly set to ground truth.
			\item The perceptual loss computes the $L_1$ distances between both $\mathbf{I}_{out}$ and $\mathbf{I}_{comp}$ and the ground truth, but after projecting these images into higher level feature spaces using an ImageNet pre-trained VGG-16.
		\end{itemize}
	\end{frame}
	
	
	\begin{frame}{Loss function}
		\centering
		\begin{itemize}
			\item The style-loss term is similar to the perceptual loss, but an autocorrelation (Gram matrix) is performed on each feature map before applying the $L_1$.
			\item As before, $\mathcal{L}_{style_{comp}}$ uses the raw output image but with non-hole pixels set to ground truth.
			\item The total variation loss $\mathcal{L}_{tv}$ represents the smoothing penalty on $R$, where $R$ is the region of 1-pixel dilation of the hole region, still evaluated using the $L_1$ norm.
		\end{itemize}
	\end{frame}
	
	
	\begin{frame}{Qualitative Results}
			\begin{columns}[c]
			
			\begin{column}{0.65\textwidth}
				\centering
				\includegraphics[height=0.8\textheight,width=\columnwidth, keepaspectratio]{images/PConvQResults1.png}
			\end{column}
			
			\begin{column}{0.5\textwidth}
				\begin{itemize}
					\item Good results for 0.1 corruption ratio, regardless of mask shape.
					\item Works well for small occlusions and can remove random noise from sap.
					\item Harder to reconstruct images in the borders.
				\end{itemize}
			\end{column}
			
		\end{columns}
	\end{frame}
	
		\begin{frame}{Qualitative Results}
			\centering
			\includegraphics[height=0.6\textheight, width=0.8\linewidth, keepaspectratio]{images/PConvQResults2.png}
		
			\begin{itemize}
				\item Frequent artifacts and colour discrepancies.
				\item When an object is mostly occluded, it struggles to generate a plausible filling, sometimes completely failing.
			\end{itemize}

	\end{frame}
	
	
	\section{Contextual Attention}
	\begin{frame}
		\begin{center}
			\Huge{\textbf{Contextual Attention}}
		\end{center}
	\end{frame}
	
	
	\begin{frame}{Goals}
		\centering
		\begin{itemize}
			\item Combine the ability of deep learning based approaches to generate visually plausible inpainting with patch synthesis methods's to borrow textures from the surrounding regions.
			\item Deal with multiple holes at arbitrary locations and with variable sizes.
			\item Integrate attention mechanism in the network: learn where to borrow or copy feature information from known background patches to generate missing patches.
			\item We implement a pre-trained model, which was trained on ImageNet for 430000 iterations.
		\end{itemize}
	\end{frame}
	
	
	\begin{frame}{Contextual Attention Layer}
		\centering
		\includegraphics[width=\textwidth, height=0.5\textheight]{images/contextual_attention_layer.jpg}
		\begin{itemize}
			\item Match features of missing pixels (foreground) to surroundings (background) via cosine similarity $\rightarrow$ weigh it with scaled softmax to get attention score for each pixel.
			\item Further encourage coherency of attention by propagation $\rightarrow$ left-right followed by a top-down propagation (kernel size k).
		\end{itemize}
	\end{frame}
	
	
	\begin{frame}{Network Architecture: $1^{st}$ stage}
		\centering
		\includegraphics[width=\textwidth, height=0.35\textheight]{images/framework.jpg}
		\begin{itemize}
			\item Two-stage coarse-to-fine network architecture where the first network makes an initial coarse prediction, and the second network takes the coarse prediction as inputs and predict refined results.
			\item The coarse network is a fully convolutional encoder-decoder architecture with dilated convolutions.
			\item It is trained with the reconstruction loss explicitly; no $BN$ layers and ELU as activation.
		\end{itemize}
	\end{frame}
	
	
	\begin{frame}{Spatially Discounted Loss}
		\centering
		\begin{itemize}
			\item Since the original image is the only ground truth to compute a reconstruction loss, strong enforcement of reconstruction loss in masked pixels may mislead the training process of the convolutional network.
			\item Missing pixels near the hole boundaries have much less ambiguity than those pixels closer to the center of the hole.
			\begin{center}
				$\downarrow$
			\end{center}
			\item Weight mask \(\mathbf{M}\) where the weight of each pixel in the mask is computed as \(\gamma^{\ell}\), where \(\ell\) is the distance of the pixel to the nearest known pixel.
		\end{itemize}
	\end{frame}
	
	
	\begin{frame}{Network Architecture: $2^{nd}$ stage}
		\centering
		\includegraphics[width=\textwidth, height=0.35\textheight]{images/unified_attention_network.jpg}
		\begin{itemize}
			\item Two parallel encoders: the bottom encoder specifically focuses on hallucinating contents with layer-by-layer (dilated) convolution, while the top one tries to attend on background features of interest.
			\item Output features from two encoders are aggregated and fed into a single decoder to obtain the final output.
			\item Color to indicate the relative location of the most interested background patch for each foreground pixel.
		\end{itemize}
	\end{frame}
	
	
	\begin{frame}{$2^{nd}$ stage Loss}
		\centering
		\begin{itemize}
			\item Reconstruction loss and WGAN-GP loss (on both global and local outputs
			of the second-stage refinement network to enforce global and local consistency).
			\item The pixel-wise reconstruction loss directly regresses
			holes to the current ground truth image, while WGANs implicitly learn to match potentially correct images and train the generator with adversarial gradients.
			\item The gradient penalty (applied only to pixels inside the holes) contributes to preventing mode collapse and balances the two critics.
			\item No perceptual, style or total variation losses used, as they do not bring noticeable improvements in this case.
			\begin{equation*}
				\begin{split}
					\min_G \max_{D \in \mathcal{D}} & \mathbb{E}_{\mathbf{x} \sim \mathbb{P}_r}[D(\mathbf{x})] - \mathbb{E}_{\tilde{\mathbf{x}} \sim \mathbb{P}_g}[D(\tilde{\mathbf{x}})] \\
					& + \lambda \mathbb{E}_{\hat{\mathbf{x}} \sim \mathbb{P}_{\hat{\mathbf{x}}}} \left[ \left( \lVert \nabla_{\hat{\mathbf{x}}} D(\hat{\mathbf{x}}) \odot (\mathbf{1} - \mathbf{m}) \rVert_2 - 1 \right)^2 \right]
				\end{split}
			\end{equation*}
		\end{itemize}
	\end{frame}
	
	
		\begin{frame}{Qualitative Results}
		\begin{columns}[c]
			
			\begin{column}{0.65\textwidth}
				\centering
				\includegraphics[height=0.8\textheight,width=\columnwidth, keepaspectratio]{images/PAttentionQResults1.png}
			\end{column}
			
			\begin{column}{0.5\textwidth}
				\begin{itemize}
					\item Very good results, even with high SAP corruption.
					\item When only a small portion of an object is missing it can reconstruct it well.
					\item Can blend well the inpainted portion with the rest of the image in patches with few colours.
				\end{itemize}
			\end{column}
			
		\end{columns}
	\end{frame}
	
	\begin{frame}{Qualitative Results}
		\centering
		\includegraphics[height=0.6\textheight, width=0.8\linewidth, keepaspectratio]{images/PAttentionQResults2.png}
		
		\begin{itemize}
			\item If the object is mostly occluded, usually the regenerated portion imitates the "likely" background, effectively editing out the item (even if with some artifacts). 
		\end{itemize}
	\end{frame}
	
		\section{Contextual Attention}
	\begin{frame}
		\begin{center}
			\Huge{\textbf{Iizuka Model}}
		\end{center}
	\end{frame}
	
	
	% ===========================================================
	% =============== IIZUKA MODEL ==============================
	% ===========================================================
	\begin{frame}{Iizuka Model: Overview}
		\begin{itemize}
			\item Two levels: \textbf{completion} + \textbf{discriminator network}
			\item Three main components:
			\begin{itemize}
				\item \textbf{Completion network}: encoder--decoder (fully convolutional)
				\item \textbf{Global discriminator}: enforces semantic consistency
				\item \textbf{Local discriminator}: enforces fine details in missing region
			\end{itemize}
			\item Encoder: strided convolutions (avoiding pooling $\rightarrow$ less blur)
			\item Decoder: deconvolutions to restore resolution
			\item Discriminators outputs combined $\rightarrow$ real vs fake classification
		\end{itemize}
		
		\centering
		\includegraphics[width=1\linewidth]{images/images_benedetto/iizuka_architecture.png}
	\end{frame}
	
	% ============= LOSS FUNCTIONS =================
	\begin{frame}{Iizuka Model: Loss Functions}
		
		\begin{itemize}
			\item \textbf{Reconstruction loss:} MSE computed only on missing region
		\end{itemize}
		
		\begin{equation*}
			L(x, M_c) = \left\| M_c \odot (C(x, M_c) - x) \right\|_2^2
		\end{equation*}
		
		where $x$ is the ground truth image, $C(x, M_c)$ the prediction, $M_c$ the binary mask, and $\odot$ element-wise multiplication.
		
		\vspace{0.3cm}
		
		\begin{itemize}
			\item \textbf{Adversarial loss:} 2-player game loss
		\end{itemize}
		
		\begin{equation*}
			\min_C \max_D \; \mathbb{E}\big[ \log D(x, M_d) + \log(1 - D(C(x, M_c), M_c)) \big]
		\end{equation*}
		
		where $M_d$ is a random mask and $D$ the joint discriminator (global + local).
		
		\vspace{0.2cm}
		By jointly minimizing and maximizing conflicting objectives, reaching stability is hard. So the training strategy is important.
		
	\end{frame}
	
	% ============== TRAINING ================
	\begin{frame}{Iizuka Model: Training Strategy}
		
		\begin{itemize}
			\item \textbf{Training} performed in three phases:
			\begin{itemize}
				\item Pretrain completion network (reconstruction loss)
				\item Train discriminators (real vs fake classification)
				\item Joint adversarial training
			\end{itemize}
			
			\item Masks: random rectangular holes (filled with dataset mean)
			
			\item \textbf{Post-processing:} Poisson blending improves boundary continuity, but may introduce color inconsistencies (it preserves gradients, not absolute intensities)
			
			\item We used the Iizuka model pre-trained on Places2 dataset ($\sim$8M images, batch size 96)
		\end{itemize}
		
	\end{frame}
	
	
	
	% ===========================================================
	% =============== IIZUKA RESULTS ============================
	% ===========================================================
	\begin{frame}{Iizuka: Qualitative Results (1)} 
		\begin{columns}[c] 
			
			\begin{column}{0.65\textwidth} 
				\centering 
				\includegraphics[height=0.85\textheight, width=0.85\linewidth, keepaspectratio]{images/images_benedetto/all_masks_comparison.png} 
			\end{column} 
			
			\begin{column}{0.6\textwidth} 
				% Fixed height container matching the image height to distribute items evenly
				\vbox to 0.6\textheight{
					\begin{itemize} 
						\item SAP mask
						\vfill
						\item Rectangular mask
						\vfill
						\item Circular mask
						\vfill
						\item Random walk mask
					\end{itemize}
				}
			\end{column} 
			
		\end{columns} 
	\end{frame}
	
	% ============== GOOD COMPARISON ================
	\begin{frame}{Iizuka: Qualitative Results (2)} 
		\begin{columns}[c] 
			
			\begin{column}{0.65\textwidth} 
				\centering 
				\includegraphics[height=1\textheight, width=1\linewidth, keepaspectratio]{images/images_benedetto/good_comparison.png} 
			\end{column} 
			
			\begin{column}{0.5\textwidth} 
				% Fixed height container matching the image height to distribute items evenly
				\vbox to 0.6\textheight{
					\begin{itemize} 
						\item A good completion
						\vfill
						\item Object shape reconstruction
						\vfill
						\item Object removal
					\end{itemize}
				}
			\end{column} 
			
		\end{columns} 
	\end{frame}
	
	% ============== BAD COMPARISON ================
	\begin{frame}{Iizuka: Qualitative Results (3)} 
		\begin{columns}[c] 
			
			\begin{column}{0.65\textwidth} 
				\centering 
				\includegraphics[height=1\textheight, width=1\linewidth, keepaspectratio]{images/images_benedetto/bad_comparison.png} 
			\end{column} 
			
			\begin{column}{0.5\textwidth} 
				% Fixed height container matching the image height to distribute items evenly
				\vbox to 0.7\textheight{
					\begin{itemize} 
						\item The presence of structured patterns make artifacts more visible to the human eye
						\vfill
						\item Color inconsistencies (bottom left corner)
						\vfill
						\item The lack of sufficient boundary constraints in the Poisson formulation can lead to unstable solutions
					\end{itemize}
				}
			\end{column} 
			
		\end{columns} 
	\end{frame}
	
	% ============== BLENDING COMPARISON 1 ================
	\begin{frame}{Iizuka: blending comparison (1)} 
		\centering 
		\includegraphics[height=1\textheight, width=1\linewidth, keepaspectratio]{images/images_benedetto/iizuka_blending_comparison_1.png} 
	\end{frame}
	
	% ============== BLENDING COMPARISON 2 ================
	\begin{frame}{Iizuka: blending comparison (2)} 
		\centering 
		\includegraphics[height=1\textheight, width=1\linewidth, keepaspectratio]{images/images_benedetto/iizuka_blending_comparison_2.png} 
	\end{frame}
	
	
	
	
	% ===========================================================
	% ==================== METRICS ==============================
	% ===========================================================
	\section{Evaluation Metrics}
	\begin{frame}
		\begin{center}
			\Huge{\textbf{Evaluation Metrics}}
		\end{center}
	\end{frame}
	
	
	\begin{frame}{Evaluation Metrics: $L_1$ and PSNR}
		
		Given a ground truth image $\mathrm{I}$ and a prediction $\hat{\mathrm{I}}$, both metrics compare them pixel-wise.
		
		\begin{itemize}
			
			\item \textbf{$L_1$ error:} Mean absolute difference over pixels. 
			Robust to outliers, but sensitive to small spatial misalignments.
			
			$$
			L_1 = \frac{1}{N} \sum_{i=1}^{N} |\mathrm{I}_i - \hat{\mathrm{I}}_i|
			$$
			
			\item \textbf{PSNR} (Peak Signal-to-Noise Ratio): Based on MSE, thus penalizes large errors more strongly. More sensitive to outliers than $L_1$. Expressed in dB. 
			
		\end{itemize}
		
		\begin{equation*}
			MSE = \frac{1}{N} \sum_{i=1}^{N} (\mathrm{I}_i - \hat{\mathrm{I}}_i)^2, 
			\qquad 
			PSNR = 10 \log_{10}\left(\frac{MAX^2}{MSE}\right)
		\end{equation*}
		
		\vspace{0.2cm}
		\textbf{Limitations}: Both are simple and easy to interpret, but do not reflect perceptual quality and are sensitive to misalignment.
		
	\end{frame}
	
	
	% ================ SSIM =====================
	\begin{frame}{Evaluation Metric: SSIM}
		
		Given two images $x$ and $y$, \textbf{SSIM (Structural Similarity Index)} measures similarity in terms of perceived structure rather than pixel-wise error. It is computed locally on small patches ($11\times 11$) and averaged over the image.
		
		\begin{equation*}
			SSIM(x,y) = \frac{(2\mu_x \mu_y + C_1)(2\sigma_{xy} + C_2)}{(\mu_x^2 + \mu_y^2 + C_1)(\sigma_x^2 + \sigma_y^2 + C_2)}
		\end{equation*}
		
		It compares luminance, contrast, and structure by merging the three terms:
		{\small
			$$
			l(x,y) = \frac{2\mu_x \mu_y + C_1}{\mu_x^2 + \mu_y^2 + C_1},
			\quad
			c(x,y) = \frac{2\sigma_x \sigma_y + C_2}{\sigma_x^2 + \sigma_y^2 + C_2}
			\quad
			s(x,y) = \frac{\sigma_{xy} + C_3}{\sigma_x \sigma_y + C_3}
			$$
		}
		
		\textbf{Limitations}: still sensitive to misalignment, favors smooth results, and ignores high-level semantics.
		
	\end{frame}
	
	
	
	
	% ================ LPIPS =====================
	\begin{frame}{Evaluation Metric: LPIPS}
		
		\textbf{LPIPS (Learned Perceptual Image Patch Similarity)} measures perceptual similarity between images $x$ and $x_0$ by comparing deep feature representations extracted from a pretrained CNN (e.g., VGG, AlexNet).
		
		\begin{equation*}
			\text{LPIPS}(x, x_0) =
			\sum_l w_l \cdot \frac{1}{H_l W_l}
			\sum_{h,w}
			\left\| \hat{\phi}_l(x)_{h,w} - \hat{\phi}_l(x_0)_{h,w} \right\|_2^2
		\end{equation*}
		
		Using weights $w_l$, it averages normalized feature differences across multiple layers of the network. 
		
		\vspace{0.2cm}
		\textbf{Limitations:} computationally expensive, but less sensitive to small misalignments.
		
		\vspace{0.2cm}
		\textbf{Common limitation (all metrics):} requires ground-truth images and cannot fully handle the one-to-many nature of inpainting.
		
	\end{frame}
	
	% ===========================================================
	% ================ GENERAL RESULTS ==========================
	% ===========================================================
	
	\section{Results Comparison}
	\begin{frame}
		\begin{center}
			\Huge{\textbf{Results Comparison}}
		\end{center}
	\end{frame}
	
	% ============ MASK RATIO ANALYSIS ==================
	\begin{frame}{Results (1)}
		
		\begin{itemize}
			\item Increasing the corrupted area makes the task more difficult
			\item For ratios $\leq 0.25$, border and non-border masks perform similarly. For higher ratio, \textbf{bordered masks perform better}. Explanation: borders contain less semantic content than image centers.
		\end{itemize}
		
		\begin{figure}
			\centering
			\includegraphics[width=0.85\linewidth]{images/images_benedetto/metrics_vs_ratio.png}
			%\caption{Metrics as a function of masked pixel ratio. Only Partial Convolution with $L_1$ is shown; other models and metrics show similar trends.}
			\label{fig:metrics_vs_ratio}
		\end{figure}
		
	\end{frame}
	
	
	% ============== METRICS vs MASKS ================
	\begin{frame}{Results (2)} 
		\begin{figure}
			\centering 
			\includegraphics[height=1\textheight, width=1\linewidth, keepaspectratio]{images/images_benedetto/metrics_vs_masks.png} 
			\caption{Mask index order: 1-9 for sap; 10-24 for box; 25-38 for circle; 39-48 for random walk}
		\end{figure}
	\end{frame}
	
	
	% ==================== MASK TYPE COMPARISON ==============================
	\begin{frame}{Effect of Mask Type}
		
		\begin{itemize}
			\item For simple masks (box and circle with small/medium ratios) most models behave similarly (except Partial Convolution).
			\item For complex masks (SAP, random walk): clearer differences emerge. \textbf{Contextual Attention}: better on sparse noise (SAP). \textbf{Iizuka + blending}: more robust on irregular structures (random walk).
		\end{itemize}
		
		\begin{figure}
			
			\centering
			\includegraphics[width=1\linewidth]{images/images_benedetto/metrics_vs_mask_method.png}
			%\caption{Model performance across different mask types. Relative rankings vary: differences are small for regular masks, but more pronounced for irregular ones.}
			\label{fig:mask_type_comparison}
		\end{figure}
		
	\end{frame}
	
	
	
	% ============== OVERALL TABLE OF METRICS ================
	\begin{frame}{Results of the metrics} 
		
		In the Table below we see the average values of each metrics for each model. The contextual attention model is the best, closely followed by the blending Iizuka.
		
		\begin{table}[htbp]
			\centering
			%\caption{Mean values for each of the considered metrics}
			\vspace{5pt}
			\label{tab:model_comparison}
			\begin{tabular}{lcccc}
				\hline
				\textbf{Model} & \textbf{L1} & \textbf{PSNR} & \textbf{SSIM} & \textbf{LPIPS} \\
				\hline
				Partial Conv.   & 0.079 & 18.2 & 0.703 & 0.318 \\
				Context. Att.   & \textbf{0.032} & \textbf{23.0} & \textbf{0.817} & \textbf{0.169} \\
				Standard Iizuka & 0.039 & 21.3 & 0.761 & 0.243 \\
				Blending Iizuka & 0.044 & 22.0 & 0.813 & 0.192 \\
				\hline
			\end{tabular}
		\end{table}
		
		\begin{table}[htbp]
			\centering
			%\caption{Mean values for each of the considered metrics}
			\vspace{5pt}
			\label{tab:model_parameters}
			\begin{tabular}{lcc}
				\hline
				\textbf{Model} & \textbf{Total trainable parameters} & \textbf{Generator only} \\
				\hline
				Partial Conv.   & 25.8M & / \\
				Context. Att.   & 9.0M & 3.6M \\
				Iizuka model    & 51.2M & 6.1M \\
				\hline
			\end{tabular}
		\end{table}
	\end{frame}
	
	
	
	
	
	% ===========================================================
	% ================ CONCLUSION ==========================
	% ===========================================================
	\begin{frame}{Conclusions}
		
		\begin{itemize}
			\item Inpainting performance strongly depends on \textbf{mask size and structure} $\rightarrow$ task is highly context-dependent.
			
			\item \textbf{Partial Convolution}: consistently worst across all metrics.
			
			\item \textbf{Best methods}: Contextual Attention and Iizuka (+ blending), with similar overall performance.
			
			\item Standard metrics ($L_1$, PSNR, SSIM, LPIPS) have \textbf{limitations}:
			\vspace{-0.4cm}
			\begin{itemize}
				\item Require ground truth
				\item Cannot capture multiple plausible solutions (one-to-many problem)
			\end{itemize}
			
			\item \textbf{Future work}:
			\begin{itemize}
				\item Larger datasets (limited here by time and computation)
				\item Human evaluation for perceptual quality
			\end{itemize}
		\end{itemize}
		
	\end{frame}
	
	
	
\end{document}